\section{Ekstern analyse}

Den eksterne analysen ser på hvilke makroforhold og bransjestrukturelle
krefter som former Ziplines handlingsrom i det norske markedet. PESTEL brukes
til å kartlegge makroomgivelsene, mens Porters Five Forces analyserer
konkurranseforholdene innenfor selve droneleveringsindustrien.

\subsection{PESTEL}

\subsubsection{Politiske faktorer}

Norges politikk for digitalisering og grønn transport tilsier at myndighetene
vil være positivt innstilt til innovativ dronelogistikk, men det regulatoriske
landskapet er komplekst. Et konkret tegn på politisk støtte er at regjeringen
i 2025 la fram Norges første stortingsmelding om droner og samtidig bevilget
50 millioner kroner til testing av null- og lavutslippsluftfart
\parencite{stortingsmelding2025}. Norge følger også EUs regelverk gjennom EØS-%
avtalen, noe som gir forutsigbarhet, men betyr samtidig at Norge ikke har
direkte påvirkningskraft på regelverksutviklingen i EU. Selskapet må innhente
BVLOS-tillatelse (Beyond Visual Line of Sight) fra Luftfartstilsynet, da
dronene opererer utenfor pilotens synsrekkevidde
\parencite{luftfartstilsynet2023}. 

\subsubsection{Økonomiske faktorer}

Det høye kostnadsnivået i Norge gjør at Zipline må vurdere nøye om
forretningsmodellen er økonomisk bærekraftig i det norske markedet. Norge har
blant Europas høyeste lønninger og driftskostnader \parencite{ssb2024}, noe
som øker investeringsbehovet betydelig. Samtidig kan droneleveranse være mer
kostnadseffektivt enn tradisjonell logistikk i distriktene, der spredt
bosetting og krevende geografi gjør vanlig distribusjon dyrere. Den høye
kjøpekraften i befolkningen, sammen med offentlig betalingsvilje og mulige
subsidier til innovativ og grønn infrastruktur, kan gi viktige inntekts- og
finansieringsmuligheter. Zipline opererer i dollar, og svingninger i NOK/USD-%
kursen skaper en valutarisiko som må håndteres. 

\subsubsection{Sosiale faktorer}

Norges befolkningsstruktur og teknologikultur gir Zipline et godt sosialt
utgangspunkt for markedsinntreden, men lokal skepsis kan utfordre aksepten.
Norges spredte befolkning med mange distriktssamfunn langt fra tradisjonell
logistikkinfrastruktur representerer et naturlig kjernesegment for
droneleveranse. Nordmenn er generelt teknologipositive og raske til å adoptere
ny teknologi, noe som senker terskelen for aksept av en ny leveringsform
\parencite{wef2024_tech_adoption}. Samtidig setter boligstrukturen rammer for
hvor praktisk gjennomførbar droneleveranse er, ettersom eneboliger i
distriktene og forstedene typisk har privat uteareal som egner seg godt til
droneleveranse. I boligblokker og leilighetskomplekser finnes det derimot
sjelden egnede private landingsområder, noe som reiser spørsmål om sikkerhet,
ansvar og identifisering av riktig mottaker. Derimot kan støy fra
droneoperasjoner over bolig- og naturområder møte motstand i lokalsamfunn,
særlig i tettbygde strøk. 

\subsubsection{Teknologiske faktorer}

Norge tilbyr et teknologisk gunstig operasjonsmiljø for dronelogistikk, men
den spesialiserte infrastrukturen for ubemannet luftfart er fortsatt under
oppbygging. Mobildekningen er blant de beste i Europa. Ved utgangen av 2024
hadde 99{,}7~\% av husstandene grunnleggende 5G-dekning, mens om lag 70~\% av
landarealet var dekket av 5G-signal \parencite{nkom2024_5g}. Høy dekning og
lav ventetid er en forutsetning for BVLOS-operasjoner, som avhenger av
kontinuerlig telemetri og sanntids posisjonsdeling mellom drone og
fjernpilot. På bransjespesifikk infrastruktur følger Norge EUs U-space-%
rammeverk, og forskrift om U-space og iAIT trådte i kraft høsten 2024
\parencite{lovdata2024_uspace}. Forskriften etablerer et digitalt
lufttrafikksystem der droner i definerte soner må dele posisjon og innhente
flygetillatelse, og testområder i Bærum og Røros er under oppstart mens de
første sertifiserte sonene ventes i 2026
\parencite{luftfartstilsynet2024_uspace}. Full utrulling vil være en
forutsetning for skalerbar dronetrafikk i tettbygde strøk, noe som innebærer
at videre teknologisk og regulatorisk modning gjenstår før bransjen kan
operere i full skala. Kombinert med offentlige satsinger på test- og
utviklingsmiljøer for ubemannet luftfart \parencite{stortingsmelding2025} gir
dette et teknologisk økosystem som favoriserer aktører som etablerer seg
tidlig. 

\subsubsection{Miljømessige faktorer}

Miljømessige forhold representerer et klart fortrinn for Zipline ved en
etablering i Norge. Elektriske droner har nullutslipp under drift, og
kombinert med at nær 100\,\% av produksjonen av elektrisitet i Norge er
fornybar \parencite{fornybarnorge}, blir droneleveranse tilnærmet
karbonnøytralt. Sammenlignet med dieseldrevne varebiler kan miljøgevinsten
derfor være betydelig. Også sammenlignet med sykkelbud tilbyr Zipline et mer
skalerbart lavutslippsalternativ over lengre distanser og i krevende terreng.
Støy kan likevel utgjøre en miljøutfordring, særlig i naturområder og
boligstrøk. Zipline forsøker imidlertid å redusere denne belastningen, som
nevnt tidligere. I tillegg viser forskning at droner kan forstyrre dyrelivet,
og særlig fugler, som kan oppfatte dronene som rovfugler og reagere med
stress, fluktatferd eller direkte angrep mot dronen \parencite{afridi2025}.
Dette gjør det viktig at Zipline planlegger flyruter som unngår hekkeområder
og sårbar natur, særlig i Norges mange verneområder og langs kjente
fugletrekkruter. Norges ambisiøse klimamål og høye miljøbevissthet i
befolkningen styrker Ziplines posisjon som et bærekraftig alternativ til
tradisjonell logistikk.

\subsubsection{Juridiske faktorer}

Det juridiske rammeverket representerer en operasjonell risiko ved en
inntreden i det norske markedet. EUs regelverk, som Norge følger gjennom EØS-%
avtalen, skiller juridisk mellom autonome og ikke-autonome droner
\parencite{euuas2021}. Ettersom Zipline benytter en fjernpilot (RPIC) med
kommando over dronen, klassifiseres den som ikke-autonom og underlegges et
mildere regelverk. Videre stiller norsk personvernlovgivning strenge krav til
innsamling, lagring og sletting av data fra dronenes kamera og sensorer
\parencite{personopplysningsloven2018}.
Ansvarsforholdet ved eventuell skade på person eller eiendom er juridisk
komplekst, og norsk produktansvarslov stiller klare krav til
erstatningsansvar. Europeisk regelverk krever dessuten at droneoperatører
tegner ansvarsforsikring \parencite{eu785_2004}, noe som kan utgjøre en
merkbar kostnad.

\subsection{Porters Five Forces}

Porters rammeverk fastslår at en aktørs lønnsomhetspotensial bestemmes av fem
konkurransekrefter, og at bransjens samlede fortjenestenivå settes av
kreftenes styrke \parencite[35]{porter1980}. Analysen under anvender
rammeverket på droneleveringsindustrien sett fra Ziplines posisjon i det
norske dagligvare- og restaurantmarkedet, og leses i sammenheng med VRIO- og
PESTEL-analysen.

\textbf{Droneleveringsindustrien i Norge kjennetegnes av høy
konkurranseintensitet fra substitutter og moderate til høye inngangsbarrierer.
Det strukturelle handlingsrommet er størst for kapitalsterke, regulatorisk
forankrede førsteetablerere, og dette vinduet er tidsbegrenset.}

\begin{table}[h!]
\centering
\caption{Oppsummering -- Porters Five Forces, droneleveringsindustrien i
Norge}
\label{tab:five_forces}
\renewcommand{\arraystretch}{1.4}
\begin{tabularx}{\textwidth}{
  >{\bfseries}p{5.2cm}
  >{\centering\arraybackslash}p{2.5cm}
  X
}
\toprule
\rowcolor{tableheader}\color{white}Kraft &
  \color{white}Styrke &
  \color{white}Nøkkeldriver \\
\midrule
\rowcolor{tablerow}
Trussel fra nye aktører    & \cellcolor{medorange} Moderat    & Høy
kapitalterskel, men attraktivt for Big Tech \\
Leverandørenes makt        & \cellcolor{medorange} Moderat    & Konsentrert
batterimarked, spesialisert hardware \\
\rowcolor{tablerow}
Kundenes makt              & \cellcolor{highred}   Høy        & Null
byttekostnad, mange substitutter \\
Trussel fra substitutter   & \cellcolor{highred}   Høy        &
Foodora/Wolt/Oda dekker allerede kjernebehovet \\
\rowcolor{tablerow}
Rivalisering i industrien  & \cellcolor{lowgreen}  Lav--moderat  & Aviant er
etablert i Norge, men med liten skala; Wing og Amazon nærmer seg \\
\bottomrule
\end{tabularx}
\end{table}

\subsubsection{Kraft 1 -- Trussel fra nye aktører
\normalfont\textit{(moderat)}}

For generelle aktører er inngangsbarrierene høye. Kapitalkrav på anslagsvis
100--300~MUSD, BVLOS-godkjenning fra Luftfartstilsynet under
EUs U-space-rammeverk \parencite{luftfartstilsynet2023} og vanskelig
repliserbare læringseffekter fra autonom operasjon \parencite[7]{porter1980}
utgjør alle betydelige hindre. Overfor ressurssterke teknologiselskaper er de
likevel ikke tilstrekkelige. Wing (Alphabet) opererer allerede kommersielt i
Finland og er dermed geografisk nær; Amazon Prime Air er FAA-godkjent og
planlegger europeisk ekspansjon \parencite{wing2024, amazonair2024}. Begge har
kapital og eksisterende retailpartnerskap som gjør rask markedsinntreden
mulig. Det avgjørende er derfor ikke barrierenes absolutte høyde, men
førsteetableringsfortrinnet.

\subsubsection{Kraft 2 -- Leverandørenes forhandlingsmakt
\normalfont\textit{(moderat)}}

Høyenergibatterier til langdistanse VTOL-operasjoner produseres av et fåtall
kinesiske og koreanske produsenter, blant dem CATL, LG Energy Solution og
Panasonic \parencite{bloombergnef2024}. Porter identifiserer slik
konsentrasjon som en kilde til sterk leverandørmakt
\parencite[27]{porter1980}. En vel så viktig driver i norsk kontekst er
maktbalansen overfor innholdspartnere. Uten attraktive dagligvare- og
restaurantpartnere har tjenesten
begrenset verdi for sluttkunden. Ettersom slikt samarbeid ennå ikke er
etablert i Norge, er det Zipline som trenger partnerne i forhandlingsfasen,
ikke omvendt.

\subsubsection{Kraft 3 -- Kundenes forhandlingsmakt \normalfont\textit{(høy)}}

Norske forbrukere har verken kontraktsmessige bindinger, teknologisk lock-in
eller nettverkseffekter som binder dem til én leveringsplattform.
Byttekostnader, som Porter identifiserer som en sentral driver av kundemakt
\parencite[24]{porter1980}, er dermed i praksis fraværende i dette
markedet. Kombinert med at betalingsviljen for sub-30-minutters dronelevering
ennå er udokumentert blant norske forbrukere \parencite[7]{pfeifer2005},
begrenser dette Ziplines prisingsevne. Resultatene fra prosjektets
spørreundersøkelse vil gi et første empirisk grunnlag for å estimere
betalingsviljen, og bør informere prismodellen direkte.

\subsubsection{Kraft 4 -- Trussel fra substitutter \normalfont\textit{(høy)}}

Porter definerer et substitutt som et produkt som fyller samme funksjon for
kjøperen \parencite[23]{porter1980}, her tilgang på mat og dagligvarer innen
akseptabel tid. Foodora, Wolt og Oda løser allerede «last mile»-problemet
\parencite{gevaers2011} til en pris markedet aksepterer, og dagligvarebutikker
innen gangavstand dekker de akutte behovene. For forbrukere i tettbygde
strøk, som utgjør en stor andel av Ziplines potensielle bykunder, er
droneleveransens verdiproposisjon primært tidsbesparelse, ettersom
tilgjengeligheten allerede er dekket av substituttene. Ziplines kombinasjon
av autonom kostnadsstruktur og kortere leveringstid gir i prinsippet et bedre
pris-ytelsesforhold enn substituttene, men substituttenes etablerte
tilgjengelighet og kundeaksept utgjør likevel et strukturelt konkurransepress
inntil dette er dokumentert i det norske markedet.

\subsubsection{Kraft 5 -- Rivalisering mellom eksisterende aktører
\normalfont\textit{(lav--moderat)}}

Det norske droneleveringsmarkedet er ikke ubeskrevet. Aviant (Trondheim,
grunnlagt 2020) opererer kommersielt i Lillehammer og Røros med BVLOS-%
sertifisering, 17~km leveringsradius og partnere som ICA og Gausdal
Landhandleri \parencite{aviant2024}. Aviant har hentet totalt 7,7~MUSD i
finansiering \parencite{cbinsights_aviant} mot Ziplines 600~MUSD i siste
runde alene \parencite{korosec2026}. Rivaliseringstrykket fra Aviant er
dermed reelt,
men foreløpig begrenset av skala. Indirekte rivalisering fra Foodora, Wolt og
Porterbuddy er høyere og mer umiddelbar, ettersom enhver vekst i Ziplines
volum primært skjer på bekostning av disse. Etter hvert som EUs U-space-%
rammeverk liberaliseres ytterligere \parencite{euuas2021}, vil også Wing og
Amazon møte lavere barrierer. Samlet gir dette en rivaliseringsintensitet som
er lav--moderat i dag og strukturelt stigende.

\subsection{Strategisk implikasjon}
Samlet peker PESTEL- og Porter-analysene mot et krevende, men tidsbegrenset
mulighetsvindu for Zipline i Norge. På makronivå er forholdene overveiende
gunstige. Norsk politikk støtter grønn luftmobilitet gjennom dedikerte
bevilgninger og regelverksutvikling \parencite{stortingsmelding2025}, det
høye norske lønnsnivået forsterker den strukturelle kostnadsfordelen ved
autonom drift \parencite{ssb2024}, og digital infrastruktur som 5G er
tilstrekkelig utbygd til å støtte BVLOS-operasjoner \parencite{nkom2024_5g}.
Bransjebildet gir en klar motvekt. Substituttene Foodora, Wolt og Oda dekker
allerede last mile-behovet til en pris markedet aksepterer, og kombinert
med tilnærmet null byttekostnad gir det norske forbrukere betydelig
forhandlingsmakt. Rivaliseringen fra Aviant er foreløpig begrenset av
skala, mens Wing og Amazon utgjør et latent inntredepress som vil tilta når
U-space- og BVLOS-regelverket modner \parencite{euuas2021,
lovdata2024_uspace}. Den strategiske implikasjonen er at Zipline bør
konsentrere seg om geografiske segmenter der substituttene fungerer dårlig,
særlig distriktsområder med spredt bosetting der rask levering har høy reell
verdi, og samtidig bruke bærekraftsprofilen aktivt overfor politiske
beslutningstakere for å sikre offentlig finansiering i etableringsfasen.
Tidsvinduet som lav lokal rivalisering og politisk velvilje åpner er
midlertidig, og førsteetablererfortrinnet avhenger av at regulatorisk
arbeid og partnerskap igangsettes tidlig. Strategiske valg konkretiseres
videre i markedsstrategikapittelet.
